\section{Raffinement des histoires et conservation de l’unité
}
\label{nuclear:1}
\subsection{Génération d’émissions et domaine de réalisation}

La construction part des émissions d’APP et de leurs prolongements enrichis. Une émission conserve les deux feuillets et leurs divisions


\[
\delta^\diamond=\operatorname{res}^\diamond+9\operatorname{quo}^\diamond,
\qquad \diamond\in\{\Sigma,\Pi\}.
\]

TRIT produit l’orientation et la retenue par

\[
\Delta(\epsilon)=\operatorname{or}(\epsilon)+3\kappa(\epsilon),
\qquad \operatorname{or}(\epsilon)\in\{-1,0,+1\}.
\]

TPK transporte la paire complète de feuillets et met à jour la mémoire par

\[
U_{\epsilon_2\epsilon_1}=U_{\epsilon_2}U_{\epsilon_1},\qquad
\mathsf U_\epsilon(m)=\mathsf C_\epsilon m+\kappa(\epsilon),\qquad
\kappa(\epsilon_2\epsilon_1)
=\kappa(\epsilon_2)+\mathsf C_{\epsilon_2}\kappa(\epsilon_1).
\]

Le cocycle rend associative la composition affine. Le bord du chemin se télescope, mais son registre ordonné est un mot ordonné, non la somme de ses entrées. L’odomètre incrémente la mémoire verticale au retour de la phase ; le retour n’identifie pas les états complets. Le transfert réduit du catalogue ne remplace pas l’action des émissions sur l’arbre ; la construction de \cref{iv:cont:historias,iv:cont:concatenacion} conserve cette distinction.

L’état enrichi conserve les cellules intermédiaires, les deux feuillets, les résidus, les quotients, l’orientation, la retenue, l’horloge, la mémoire, la frontière, la route et le registre ordonné. Une émission admissible met conjointement à jour ces champs. La troncature efface exactement sa dernière contribution. Un prolongement neutre donne une section de la troncature, mais enregistre une étape : ce n’est pas l’absence d’histoire. De là proviennent la finitude, la non-vacuité et la surjectivité de tous les niveaux, puis la survie totale.

Cette partie du développement porte sur la réalisation analytique de ce système et sur sa mémoire. Aucune valeur numérique cible ni constante conventionnelle n’intervient comme entrée. Les coefficients probabilistes suivants proviennent uniquement du nombre d’émissions sortantes effectives.


\subsection{Raffinement sur la multisection complète
}

Fixons un préfixe enrichi réel \(x\) de profondeur \(k\). Nous écrivons


\[
H_n=\operatorname{Term}_{k,k+n}(x),\qquad
\rho_n:H_{n+1}\longrightarrow H_n,
\qquad \Omega_x=\varprojlim_n(H_n,\rho_n).
\]

Chaque élément de \(H_n\) est un prolongement complet jusqu’à cet horizon, avec ses témoins d’émission. Pour \(h\in H_n\), les successeurs sont \(h\epsilon\), pour toutes les émissions \(\epsilon\in\operatorname{Out}(h)\). Les multiplicités des opérations sont conservées, même lorsque des données terminales réduites coïncident. Par \cref{iv:cont:medida} :

\[
p_h(\epsilon)=\frac1{d(h)}>0,\qquad
d(h)=\#\operatorname{Out}(h),\qquad
\mu_{n+1}(h\epsilon)=\mu_n(h)p_h(\epsilon).
\]

Par conséquent,

\begin{equation}
\sum_{\epsilon\in\operatorname{Out}(h)}\mu_{n+1}(h\epsilon)=\mu_n(h),
\qquad (\rho_n)_*\mu_{n+1}=\mu_n.
\label{nuc01:eq:1}
\end{equation}

La proposition~\ref{iv:cont:medida} construit l’unique probabilité \(\mu_x\) sur \(\Omega_x\) ayant ces masses de cylindre ; tout cylindre non vide a une masse positive. La mesure ne sélectionne pas une histoire et n’impose pas de filtre postérieur à l’existence des cinq constructions.

\begin{proposition}[Le raffinement des observables est isométrique et unital]
\label{nuc01:pullback}

Soient \(\mathscr L_n=L^2(H_n,\mu_n)\) et

\[
R_n:\mathscr L_n\longrightarrow\mathscr L_{n+1},\qquad
(R_nf)(h\epsilon)=f(h).
\]

Alors \(R_n^*R_n=I\), \(R_n1=1\), et son adjoint est

\begin{equation}
(E_ng)(h)=\sum_{\epsilon\in\operatorname{Out}(h)}p_h(\epsilon)g(h\epsilon).
\label{nuc01:eq:2}
\end{equation}

En particulier, \(E_nR_n=I\), \(E_n1=1\), et \(R_nE_n\) est la projection orthogonale sur les observables qui ne dépendent que du préfixe.
\end{proposition}

\begin{proof}
En appliquant \eqref{nuc01:eq:1},


\[
\|R_nf\|_{n+1}^2
=\sum_h\mu_n(h)|f(h)|^2\sum_\epsilon p_h(\epsilon)
=\|f\|_n^2.
\]

Le même regroupement dans le produit scalaire donne \eqref{nuc01:eq:2}. Les affirmations sur l’unité sont ponctuelles ; l’affirmation sur la projection découle de \(E_n=R_n^*\) et de \(E_nR_n=I\). De plus, \(R_n(fg)=R_nf\,R_ng\) et \(R_n\bar f=\overline{R_nf}\) : sur les observables finies, c’est aussi un homomorphisme unital d’algèbres.

\end{proof}

Par itération, \(R_{m,n}=R_{m-1}\cdots R_n\) coïncide avec le tiré en arrière par la troncature composée. Les tirés en arrière cylindriques \(R_{\infty,n}\) réalisent la limite inductive de ces espaces de Hilbert dans \(L^2(\Omega_x,\mu_x)\). Leur réunion est dense : les cylindres forment une algèbre qui engendre la tribu, et les fonctions simples de cette algèbre sont denses dans \(L^2\). Il n’est pas nécessaire de choisir une trajectoire pour construire la limite.

\subsubsection{La même isométrie dans la base des histoires}

L’identification unitaire \(W_nf(h)=\sqrt{\mu_n(h)}f(h)\) envoie \(\mathscr L_n\) sur \(\ell^2(H_n)\). En conjuguant \(R_n\), on obtient

\begin{equation}
\widehat R_n\delta_h
=\sum_{\epsilon\in\operatorname{Out}(h)}\sqrt{p_h(\epsilon)}\,\delta_{h\epsilon}.
\label{nuc01:eq:3}
\end{equation}

Les vecteurs de droite ont des supports disjoints pour des prédécesseurs distincts, parce que \(\rho_n(h\epsilon)=h\). Chaque colonne a une norme un par normalisation. Cela démontre directement l’isométrie sur les histoires enrichies déjà engendrées, sans imposer qu’une projection numérique conserve toute la mémoire. Le vecteur unité se transforme en

\[
\Omega_n=\sum_{h\in H_n}\sqrt{\mu_n(h)}\,\delta_h,
\qquad \widehat R_n\Omega_n=\Omega_{n+1}.
\]

On n’ajoute pas une copie artificielle de \(h\) à un résultat pour le récupérer : \(h\epsilon\) est précisément le prolongement que le domaine APP–TRIT–TPK a déjà produit.


\subsection{Prolongement corrélatif des cinq opérations}

La construction de \cref{iv:cont:mapa-correlativo} définit un unique registre \(\mathfrak e_\epsilon\) par émission. Ses opérations spectrale, solénoïdale, de jauge, cohomologique et déterminantale dépendent de cette même émission, avec les mêmes source, destination, orientation, retenue, frontière, route et registre ordonné. L’action cellulaire satisfait


\[
\partial\operatorname{Cell}(\epsilon)=\operatorname{Cell}(\epsilon)\partial,
\qquad
\operatorname{Cell}(\epsilon_2\epsilon_1)
=\operatorname{Cell}(\epsilon_2)\operatorname{Cell}(\epsilon_1).
\]

Le complexe et sa droite déterminante sont ceux du même objet ; le remplissage construit vérifie \(\partial h_*=z_--z_+\). Le caractère total du constructeur n’exige pas qu’une trivialisation spécialisée soit acyclique ou inversible. L’existence conjointe est ainsi conservée avant tout lecteur spécialisé.

La naturalité de \cref{iv:rectangular-natural,iv:cont:memoria-natural} s’exprime conjointement par

\begin{equation}
u_{(\ell,N'),(k,N)}\mathfrak D_{\ell,N'}(\widehat x_\ell)
=\mathfrak D_{k,N}(\tau_{\ell,k}\widehat x_\ell).
\label{nuc01:eq:4}
\end{equation}

Sa preuve n’identifie pas une arête fine isolée à une arête grossière fictive : elle restreint l’histoire des émissions et conserve les actions de leurs entrées. Le transport se compose sur les mêmes probabilités cylindriques ; les bilans utilisent l’espérance conditionnelle et son correcteur ; l’incidence et les ports utilisent leurs applications cellulaires ; l’homologie et le déterminant utilisent la même application de complexes. Ce ne sont pas cinq limites d’histoires choisies séparément.

Il y a deux directions à distinguer. Les applications d’état et de coefficients vont du fin au grossier. Les images réciproques \(R_n\) de la proposition~\ref{nuc01:pullback} vont des observables grossières aux fines. L’espérance \(E_n\) revient dans la direction de réduction. Cette distinction permet de réaliser \eqref{nuc01:eq:4} analytiquement sans inverser ses domaines.


La construction terminale de \cref{iv:cont:terminal,iv:cont:ensamblaje,iv:cont:determinantes} prolonge les terminaux complets et conserve simultanément le télescopage polaire, la matrice de Gram et les pertes, l’incidence avec mesure, le cône de retour et les données Smith–Bockstein. La promotion de l’une de ces actions au rang d’opérateur borné requiert la norme et la borne correspondante (voir le passage à la limite démontré à la fin de cette section). Une collection finie compatible ne devient pas, par son seul nom, un opérateur borné de la limite.

\subsection{Fibres de mémoire : critère de Gram et conservation de l’unité}

Pour une réalisation non scalaire, attribuons à chaque préfixe \(h\) une fibre de Hilbert \(\mathscr M_h\), et à son émission un transport linéaire typé
\(I_\epsilon:\mathscr M_h\to\mathscr M_{h\epsilon}\). Cet opérateur linéaire n’est pas identifié à la mise à jour affine \(\mathsf U_\epsilon\) sans déclarer la représentation qui la réalise.


Sur \(\mathscr H_n=\bigoplus_{h\in H_n}\mathscr M_h\), on définit

\[
(V_n\xi)_{h\epsilon}=\sqrt{p_h(\epsilon)}I_\epsilon\xi_h.
\]

La disjonction des successeurs démontre l’identité exacte de \cref{iv:cont:gram} :

\begin{equation}
\left.V_n^*V_n\right|_{\mathscr M_h}
=\sum_{\epsilon\in\operatorname{Out}(h)}p_h(\epsilon)I_\epsilon^*I_\epsilon.
\label{nuc01:eq:5}
\end{equation}

Le critère nécessaire et suffisant pour que \(V_n\) soit isométrique est que la somme de \eqref{nuc01:eq:5} soit \(I\) en chaque prédécesseur. Que chaque \(I_\epsilon\) soit isométrique est suffisant. C’est également nécessaire si, au préalable, tous les \(I_\epsilon\) sont des contractions : alors chaque \(I-I_\epsilon^*I_\epsilon\) est positif, et une somme à poids strictement positifs ne peut s’annuler que si chaque terme s’annule. Le critère de Gram caractérise l’isométrie de l’application totale ; l’isométrie de chaque transport constitue une condition suffisante, qui s’avère également nécessaire dans la classe des contractions.

Dans les histoires scalaires de \eqref{nuc01:eq:3}, \(I_\epsilon=1\), de sorte que la condition est démontrée. Dans les fibres analytiques supplémentaires, on conserve \eqref{nuc01:eq:5} et on identifie l’hypothèse concrète permettant de passer à l’isométrie.

Si les fibres possèdent des vecteurs unitaires distingués \(u_h\) et \(I_\epsilon u_h=u_{h\epsilon}\), alors


\begin{equation}
V_n\Bigl(\bigoplus_h\sqrt{\mu_n(h)}u_h\Bigr)
=\bigoplus_{h\epsilon}\sqrt{\mu_{n+1}(h\epsilon)}u_{h\epsilon}.
\label{nuc01:eq:6}
\end{equation}

L’isométrie seule n’affirme pas \eqref{nuc01:eq:6} : conserver la norme et conserver un vecteur unité distingué sont des propriétés différentes.


\subsection{Symétries et lecteurs sur le domaine commun}

Une symétrie réversible cohérente de l’arbre transporte toutes les émissions et leurs multiplicités. L’invariance des multiplicités dans \cref{iv:cont:medida} prouve qu’elle conserve les degrés et les probabilités des cylindres. Elle peut envoyer \(\Omega_x\) sur \(\Omega_{gx}\) ; elle n’agit à l’intérieur d’une même fibre de prolongements que si elle fixe le préfixe.

Par conséquent

\[
(K_gf)(g\omega)=f(\omega)
\]

définit une application unitaire \(L^2(\Omega_x,\mu_x)\to L^2(\Omega_{gx},\mu_{gx})\). La naturalité de la troncature implique \(K_{g,n+1}R_n=R'_nK_{g,n}\), et, par adjonction, l’égalité correspondante des espérances. Dans les fibres de mémoire, on exige en outre le carré


\[
S_{g,h\epsilon}I_\epsilon=I_{g\epsilon}S_{g,h}
\]

avec \(S_{g,h}\) unitaires. C’est l’opérateur d’entrelacement qui transforme une symétrie de l’arbre en symétrie de cette représentation de mémoire ; on ne l’attribue pas nominalement à tout groupe exceptionnel ultérieur.


Pour la double lecture, on prend la jauge initiale et son transport préfixal déclaré dans \cref{exc:doble-lectura,exc:limite-inverso}. On travaille sur une multisection commune \(\Omega_x\) où les lecteurs sont définis ; pour le lecteur archimédien, cela inclut les intervalles compacts emboîtés et la contraction des diamètres déclarés dans \eqref{eq:cilindro}. C’est une prémisse sur la réalisation du lecteur, et non un nouveau critère d’existence du constructeur conjoint. Les identités finies sont

\begin{equation}
r_{n,m}Q_n=Q_m\rho_{n,m},\qquad
a_{n,m}R^{\rm arq}_n=R^{\rm arq}_m\rho_{n,m}.
\label{nuc01:eq:7}
\end{equation}

Ils déterminent la lecture conjointe à la limite. Si \(L=(\operatorname{ev}_{\rm arq},Q_\infty)\) et \(\nu=L_*\mu_x\), alors

\[
C_L:L^2(Y,\nu)\longrightarrow L^2(\Omega_x,\mu_x),
\qquad C_Lf=f\circ L,
\]

est une isométrie, car l’égalité des normes est la définition de la mesure image. Elle conserve l’unité et ses images sont les observables mesurables par rapport à l’information exprimée par \(L\). Pour des lecteurs finis compatibles, \eqref{nuc01:eq:7} implique également la cohérence des mesures images.

Cette isométrie n’équivaut pas à ce que le lecteur distingue tous les états. La construction de \cref{exc:limite-inverso} distingue le fait qu’un même courant visible avec une mémoire cachée différente peut avoir la même incidence. L’adjoint de \(C_L\) est l’espérance par rapport au lecteur ; \(C_LC_L^*\) conserve exactement cette partie observable. Si une réalisation élargie \((L,M)\) est injective, borélienne et d’inverse mesurable sur son image — par exemple, une injection continue de l’espace compact des histoires dans un espace de Hausdorff approprié —, son image réciproque identifie effectivement tout \(L^2\). La mémoire requise est celle déjà conservée dans l’état, non une nouvelle sélection rétrospective.


\subsection{Théorème de passage à la limite pour la charge observable/mémoire}

Cette section type le lien avec une identité de conservation de mémoire. Le raffinement canonique des observables fournit les applications isométriques nécessaires. Une réalisation supplémentaire dans les fibres de mémoire utilise le critère de Gram de \eqref{nuc01:eq:5}.

Soient deux systèmes inductifs d’espaces de Hilbert \((\mathscr H_n,R_n)\) et \((\mathscr K_n,S_n)\), à inclusions isométriques. Pour chaque \(n\), soient

\[
J_n:\mathscr H_n\to\mathscr K_n\text{ isométrique},\quad
P_n=J_nJ_n^*,\quad U_n\text{ unitaire sur }\mathscr K_n,
\]

et \(\widetilde A_n\) autoadjoint borné, avec \(\sup_n\|\widetilde A_n\|<\infty\). Nous définissons

\[
A_n=J_n^*\widetilde A_nJ_n,\quad
T_n=J_n^*U_nJ_n,\quad
\eta_n=(I-P_n)U_nJ_n.
\]

Supposons la conservation et la séparation à chaque niveau,

\begin{equation}
U_n^*\widetilde A_nU_n=\widetilde A_n,
\qquad [\widetilde A_n,P_n]=0,
\label{nuc01:eq:8}
\end{equation}

ainsi que les carrés de raffinement suivants :

\begin{equation}
J_{n+1}R_n=S_nJ_n,\qquad
P_{n+1}S_n=S_nP_n,
\label{nuc01:eq:9}
\end{equation}

\begin{equation}
U_{n+1}S_n=S_nU_n,\qquad
\widetilde A_{n+1}S_n=S_n\widetilde A_n.
\label{nuc01:eq:10}
\end{equation}

\begin{theorem}
Ces collections induisent des opérateurs bornés dans les limites inductives, avec \(J_\infty\) isométrique et \(U_\infty\) unitaire, et conservent

\begin{equation}
A_\infty
=T_\infty^*A_\infty T_\infty
 +\eta_\infty^*\widetilde A_\infty\eta_\infty.
\label{nuc01:eq:11}
\end{equation}

La même identité vaut à chaque profondeur finie. Si \(\widetilde A_n\geq0\), la charge retenue est positive. Sans positivité, on conserve une forme autoadjointe, pas nécessairement une énergie positive.
\end{theorem}

\begin{proof}
Nous décomposons \(U_nJ_n=J_nT_n+\eta_n\). Les deux termes appartiennent à \(P_n\mathscr K_n\) et à \((I-P_n)\mathscr K_n\). La commutation de \eqref{nuc01:eq:8} annule les termes croisés de la forme \(\widetilde A_n\). En appliquant l’autre identité de \eqref{nuc01:eq:8},


\[
A_n=J_n^*U_n^*\widetilde A_nU_nJ_n
=T_n^*A_nT_n+\eta_n^*\widetilde A_n\eta_n.
\]

La seconde égalité de \eqref{nuc01:eq:9}, jointe à la première, implique
\(J_{n+1}^*S_n=R_nJ_n^*\) : il suffit d’écrire
\(P_{n+1}S_n=S_nJ_nJ_n^*=J_{n+1}R_nJ_n^*\) et de multiplier par \(J_{n+1}^*\).
Par \eqref{nuc01:eq:9}–\eqref{nuc01:eq:10},

\[
T_{n+1}R_n=R_nT_n,\qquad
\eta_{n+1}R_n=S_n\eta_n,\qquad
A_{n+1}R_n=R_nA_n.
\]

Les opérateurs sont ainsi définis sur les réunions algébriques des systèmes. Les isométries et la borne uniforme permettent de les prolonger par continuité aux complétions. La compatibilité de \(U_n^*\) se déduit aussi de celle de \(U_n\), car les deux niveaux sont unitaires ; la limite est, par conséquent, unitaire et non seulement isométrique. Les projections \(P_n\) induisent une projection dont l’image est l’adhérence de la réunion des images de \(J_n\) ; c’est \(J_\infty J_\infty^*\). Toutes les identités persistent d’abord sur les vecteurs de niveau fini, puis par densité. Cela prouve \eqref{nuc01:eq:11}.
\end{proof}

La seule compatibilité de \(J_n\) ne remplace pas la seconde égalité de \eqref{nuc01:eq:9}. Un raffinement qui conserve l’inclusion observable mais mélange son complément avec elle peut modifier le défaut de mémoire.

\subsubsection{Réalisation nonadique explicite compatible
}

Supposons \(C_n\) unitaire dans \(\mathscr H_n\), naturel par rapport à \(R_n\), et \(A_n\) autoadjoint borné, uniformément borné, naturel et tel que \([A_n,C_n]=0\). Sur \(\mathscr K_n=\mathscr H_n^{\oplus9}\), nous définissons


\[
J_nv=\tfrac13(v,\ldots,v),\quad
S_n=R_n^{\oplus9},\quad
U_n(v_1,\ldots,v_9)=(C_nv_9,v_1,\ldots,v_8),
\quad \widetilde A_n=I_9\otimes A_n.
\]

Le facteur \(1/3\) normalise les neuf composantes. Les formules vérifient directement \eqref{nuc01:eq:8}–\eqref{nuc01:eq:10}, y compris la projection. La compression est \(T_n=(8I+C_n)/9\), et

\begin{equation}
\eta_n^*\widetilde A_n\eta_n
=A_n-T_n^*A_nT_n
=\frac8{81}(I-C_n)^*A_n(I-C_n).
\label{nuc01:eq:12}
\end{equation}

La dernière égalité utilise \(C_n^*A_nC_n=A_n\). D’après le théorème, \eqref{nuc01:eq:12} vaut aussi à la limite. Le raffinement nonadique ou le raffinement complet des histoires transporte la charge et sa mémoire lorsque \(C_n\) et \(A_n\) satisfont les carrés indiqués.


Pour \(A=I\), le télescopage fini est

\[
\|v\|^2=\|T^Nv\|^2+\sum_{j=0}^{N-1}\|\eta T^jv\|^2.
\]

Par conséquent, \(v\mapsto(\eta v,\eta Tv,\ldots,\eta T^{N-1}v,T^Nv)\) est une isométrie. La formule de charge correspondante remplace les normes par les formes de \(A\) et \(\widetilde A\). On n’élimine pas le terme terminal sans prouver sa convergence vers zéro ; la compatibilité de profondeur ne prouve pas à elle seule cette stabilité temporelle.

Si \(D_{N,n}\) désigne cette isométrie à la profondeur \(n\) et à l’horizon temporel \(N\), les carrés précédents prouvent en outre


\[
(S_n^{\oplus N}\oplus R_n)D_{N,n}
=D_{N,n+1}R_n.
\]

En effet, chaque composante satisfait \(\eta_{n+1}T_{n+1}^jR_n=S_n\eta_nT_n^j\), et la composante terminale satisfait \(T_{n+1}^NR_n=R_nT_n^N\). L’accumulation de mémoire à horizon fini et le raffinement de résolution commutent. Le même calcul prouve l’équivariance pour une symétrie représentée par des unitaires qui entrelacent \(J,U,P,\widetilde A\) ; il n’étend pas l’action aux symétries pour lesquelles ces carrés n’ont pas été construits.
